\chapter{Dementia and Driving}
\index{Dementia and Driving}

\section*{Definition and Overview}
Driving is a complex task that requires memory, attention, judgement, quick reactions, and visual skills, all of which can be affected by dementia. As dementia progresses, driving becomes unsafe, and the decision about when to stop is one of the most difficult and important issues facing people with dementia and their families. In many countries, people with dementia are legally required to notify their driver licensing authority, and medical fitness to drive must be assessed. This chapter explains the legal obligations, how driving ability is assessed, and how to manage the transition to stopping driving.

\section*{Epidemiology}
Older drivers are a growing proportion of road users, and many people with early dementia continue to drive. Research shows that people with dementia have a significantly increased risk of crashes compared with people of the same age without dementia, particularly as the condition progresses. However, in the early stages, many people with mild cognitive impairment or early dementia still drive safely. In many countries, doctors are required to report conditions that may impair driving ability to the driver licensing authority, balancing patient autonomy with public safety.

\section*{Aetiology and Pathophysiology}
\begin{itemize}
    \item \textbf{Cognitive decline:} memory loss, impaired attention, and poor judgement affect decision-making on the road
    \item \textbf{Visuospatial problems:} difficulty judging distances and processing visual information increases crash risk
    \item \textbf{Slowed reactions:} processing speed and reaction time decline
    \item \textbf{Impaired insight:} a hallmark of dementia is reduced awareness of one's own limitations, making self-assessment unreliable
    \item \textbf{Getting lost:} disorientation and memory failure cause wandering and wrong-way driving
    \item \textbf{Other symptoms:} apraxia, confusion under stress, and sleep disturbance all impair driving safety
\end{itemize}

\section*{Clinical Features}
\begin{itemize}
    \item Getting lost on familiar routes
    \item Forgetting how to operate the car or misusing pedals
    \item Slow responses, near-misses, and difficulty merging or judging gaps
    \item Confusion at intersections and in unfamiliar areas
    \item Dents and scratches on the car from minor collisions
    \item Reduced awareness of road signs and hazards
    \item Difficulty following conversations while driving or becoming easily distracted
    \item Angry reactions or agitation when others comment on driving
\end{itemize}

\section*{Differential Diagnosis}
\begin{itemize}
    \item Normal ageing changes, which are mild and may not affect driving
    \item Medication effects causing drowsiness or confusion
    \item Visual impairment, including cataracts and macular degeneration
    \item Other neurological conditions such as Parkinson's disease affecting coordination
    \item Sleep disorders causing daytime drowsiness
    \item Acute illness or delirium, which temporarily affects driving ability
\end{itemize}

\section*{When to Seek Medical Care}
A diagnosis of dementia should be discussed with a doctor in relation to driving, and the driver licensing authority should be notified. Medical review is needed if there is any concern about driving safety, after a near-miss or crash, or if the condition progresses. Family members who are worried about a relative's driving should raise their concerns with the doctor.

\textbf{Contact your local emergency services for:}
\begin{itemize}
    \item A person with dementia who is missing or has driven off and is lost (immediate danger)
    \item A motor vehicle accident with injuries
    \item Sudden confusion or collapse while driving or shortly after
\end{itemize}

\section*{Diagnosis and Investigations}
\begin{itemize}
    \item \textbf{Medical assessment:} review of the dementia diagnosis, stage, and symptoms relevant to driving
    \item \textbf{Cognitive testing:} screening tests provide information but do not alone determine fitness to drive
    \item \textbf{On-road assessment:} a practical driving test with an occupational therapist or approved assessor is the most reliable measure of driving ability
    \item \textbf{Off-road assessment:} computer-based and simulator tests of attention, reaction time, and visuospatial skills
    \item \textbf{Formal notification:} in many countries, doctors complete medical fitness-to-drive reports to the local licensing authority
    \item \textbf{Regular reassessment:} driving ability is reviewed as dementia progresses, since deterioration is expected
\end{itemize}

\section*{Management}
\begin{itemize}
    \item \textbf{Legal compliance:} notify the driver licensing authority of the diagnosis, as required by law in many countries
    \item \textbf{Fitness-to-drive review:} the licensing authority may require a medical report and practical driving assessment
    \item \textbf{Driving restrictions:} some people may be allowed to drive with conditions, such as daylight-only driving or limits on distance and familiar areas
    \item \textbf{Graduated cessation:} planning to stop driving gradually, with alternatives arranged in advance
    \item \textbf{Alternative transport:} public transport, community transport services, family support, and taxis
    \item \textbf{Emotional support:} stopping driving is a major loss of independence; counselling and support groups help with adjustment
    \item \textbf{Removing access:} keys, car keys, and car access may need to be managed when the person cannot safely drive
    \item \textbf{Vehicle adaptations:} in some cases, occupational therapy assessment can recommend modifications
\end{itemize}

\section*{Complications}
\begin{itemize}
    \item Motor vehicle crashes with injury or death to the person with dementia or others
    \item Legal consequences of driving with an un-notified condition
    \item Loss of independence, social isolation, and depression when driving stops
    \item Family conflict over the decision to stop driving
    \item Wandering and getting lost, which can expose the person to danger
    \item Financial and practical burden on families providing transport
\end{itemize}

\section*{Prognosis and Outcomes}
The ability to drive safely declines as dementia progresses, and ultimately all people with dementia must stop driving. In the early stages, some people can continue to drive with regular review and appropriate restrictions. Planning for the transition, discussing it early, and arranging alternatives make stopping driving easier and safer. While the loss of driving is difficult, most people and families adapt with support, and the safety benefits for the person and the community are substantial.

\section*{Prevention and Self-Care}
\begin{itemize}
    \item Discuss driving with the doctor at diagnosis and at each review
    \item Notify the driver licensing authority promptly
    \item Take an on-road driving assessment to get an objective opinion
    \item Plan alternative transport arrangements before they are needed
    \item Limit driving to familiar areas, daylight, and short distances if recommended
    \item Family members: raise concerns early, share responsibility, and seek support
    \item Use community transport, family rosters, and mobility services to maintain independence after stopping
    \item Keep the car keys and car access managed once driving has stopped
\end{itemize}

\section*{Sources and Further Reading}
The following reputable sources provide further information for healthcare students and patients seeking reliable, current advice about dementia and driving.

\begin{itemize}
    \item Dementia Australia. \textit{Driving and dementia: information for families.} https://www.dementia.org.au/
    \item Austroads. \textit{Assessing fitness to drive: guidelines for health professionals.} https://austroads.com.au/
    \item Transport for NSW and state licensing authorities. \textit{Medical fitness to drive requirements.} https://www.nsw.gov.au/
    \item Alzheimer's Association. \textit{Dementia and driving: when to stop.} https://www.alz.org/
    \item NHS. \textit{Dementia and driving: rules and support.} https://www.nhs.uk/conditions/dementia/
    \item Occupational Therapy Australia. \textit{Driving assessment services.} https://otaus.com.au/
\end{itemize}
